\documentclass[12pt, a4paper]{article}

% --- UNIVERSAL PREAMBLE BLOCK ---
\usepackage[top=2.5cm, bottom=2.5cm, left=3cm, right=2.5cm, headheight=15pt]{geometry}
\usepackage[T1]{fontenc}
\usepackage[utf8]{inputenc}
\usepackage[brazilian]{babel}

% Use a high-quality Sans Serif font compatible with pdflatex
\usepackage[default, osf]{sourcesanspro}
\usepackage{zi4} % Monospace font (Consolas-like)

% Add because main language is not English
\usepackage{enumitem}
\setlist[itemize]{label=-}

% --- Cores ---
\usepackage[dvipsnames, table]{xcolor}
\definecolor{pixeon}{RGB}{0, 102, 179}      % azul institucional Pixeon
\definecolor{codebg}{RGB}{245, 247, 250}
\definecolor{coderule}{RGB}{200, 210, 225}
\definecolor{gray80}{RGB}{80, 80, 80}

% --- Tipografia e Layout ---
\usepackage[expansion=false]{microtype}
\usepackage{setspace}
\onehalfspacing
\setlength{\parindent}{1.25cm}
\setlength{\parskip}{4pt}

% --- Cabeçalho / rodapé ---
\usepackage{fancyhdr}
\pagestyle{fancy}
\fancyhf{}
\fancyhead[L]{\small\color{gray80} Pixeon Image Viewer --- Relatório Técnico}
\fancyhead[R]{\small\color{gray80} Vitor N. M. da Silva}
\fancyfoot[C]{\small\color{gray80} \thepage}
\renewcommand{\headrulewidth}{0.4pt}

% --- Seções com cor ---
\usepackage{titlesec}
\titleformat{\section}
  {\large\bfseries\color{pixeon}}
  {\thesection.}{0.6em}{}[\vspace{-2pt}\color{pixeon}\rule{\linewidth}{0.6pt}]
\titleformat{\subsection}
  {\normalsize\bfseries\color{gray80}}
  {\thesubsection.}{0.5em}{}

% --- Código-fonte ---
\usepackage{listings}
\lstset{
  backgroundcolor=\color{codebg},
  basicstyle=\ttfamily\footnotesize,
  breaklines=true,
  frame=single,
  rulecolor=\color{coderule},
  framesep=4pt,
  numbers=left,
  numberstyle=\tiny\color{gray80},
  numbersep=8pt,
  keywordstyle=\color{pixeon}\bfseries,
  commentstyle=\color{OliveGreen}\itshape,
  stringstyle=\color{BrickRed},
  showstringspaces=false,
  tabsize=4,
  captionpos=b,
  language=C++
}

% --- Tabelas e Matemática ---
\usepackage{booktabs}
\usepackage{tabularx}
\usepackage{array}
\usepackage{amsmath}

% --- Caixas (tcolorbox) ---
\usepackage{tcolorbox}
\tcbuselibrary{skins, breakable}

\newtcolorbox{infobox}[1][]{
  enhanced, breakable,
  colback=codebg, colframe=pixeon,
  fonttitle=\bfseries\small,
  left=6pt, right=6pt, top=4pt, bottom=4pt,
  #1
}

% --- Links e PDF metadata (deve ser o último pacote) ---
\usepackage[
  colorlinks=true,
  linkcolor=pixeon,
  urlcolor=pixeon,
  citecolor=pixeon,
  pdftitle={Relatório Técnico — Pixeon Image Viewer},
  pdfauthor={Vitor Neuenschwander Mendes da Silva},
  pdfsubject={Desafio Técnico — Analista Desenvolvedor C++},
  pdfkeywords={Qt, C++17, Processamento de Imagem, MVC, PACS}
]{hyperref}

% ========================================================================
\begin{document}

% --- Capa ---
\begin{titlepage}
  \centering
  \vspace*{2cm}

  {\color{pixeon}\rule{\linewidth}{2pt}}
  \vspace{0.4cm}

  {\Huge\bfseries\color{pixeon} Pixeon Image Viewer}\\[0.3cm]
  {\Large\color{gray80} Medical Edition}

  \vspace{0.3cm}
  {\color{pixeon}\rule{\linewidth}{2pt}}

  \vspace{1.5cm}
  {\large\bfseries Relatório Técnico}\\[0.2cm]
  {\normalsize Desafio Técnico --- Analista Desenvolvedor C++ Júnior}\\
  {\normalsize \textbf{Código fonte: }\href{https://github.com/neuxxkk/pixeon-image-viewer}{ neuxxkk/pixeon-image-viewer} }

  \vfill

  \begin{tabular}{rl}
    \textbf{Candidato:}  & Vitor Neuenschwander Mendes da Silva \\[4pt]
    \textbf{Contato:}    & \href{mailto:vitornms@gmail.com}{vitornms@gmail.com} \\[4pt]
    \textbf{LinkedIn:}   & \href{https://linkedin.com/in/vitornms}{linkedin.com/in/vitornms} \\[4pt]
    \textbf{GitHub:}     & \href{https://github.com/neuxxkk}{github.com/neuxxkk} \\[4pt]
    \textbf{Data:}       & \today
  \end{tabular}

  \vspace{1.5cm}
  {\color{pixeon}\rule{\linewidth}{0.8pt}}
  \vspace{0.2cm}
  {\small\color{gray80}
    Tecnologias: C++17 \textbullet{} Qt 6 \textbullet{} CMake \textbullet{} QTest
  }
\end{titlepage}

% --- Sumário ---
\tableofcontents
\newpage

% ========================================================================
\section{Introdução}

Este documento descreve as decisões técnicas, metodologia e resultados obtidos
no desenvolvimento do \textbf{Pixeon Image Viewer}, aplicação desktop de
visualização e manipulação de imagens médicas desenvolvida como desafio técnico
para a vaga de Analista Desenvolvedor C++ Júnior na Pixeon.

O projeto evoluiu intencionalmente além dos requisitos mínimos do enunciado,
incorporando funcionalidades de apoio ao diagnóstico médico, como calibração de
medições em milímetros e inversão de polaridade radiográfica, refletindo o
contexto do produto da empresa.

\begin{infobox}[title={Repositório e Executável}]
  \textbf{Código-fonte:}
  \url{https://github.com/neuxxkk/pixeon-image-viewer}\\
  O repositório contém o código-fonte completo, o arquivo de projeto CMake e
  o executável Windows empacotado com \texttt{windeployqt}.
\end{infobox}

% ========================================================================
\section{Metodologia e Arquitetura}

\subsection{Padrão MVC}

O projeto adota o padrão \textbf{Model-View-Controller (MVC)} como espinha
dorsal da arquitetura, separando claramente as três responsabilidades:

\begin{table}[h!]
\centering
\caption{Distribuição de responsabilidades no padrão MVC}
\vspace{4pt}
\begin{tabularx}{\linewidth}{>{\bfseries}lX}
\toprule
Camada & Responsabilidade \\
\midrule
Model 
\\
\texttt{ImageModel}, \texttt{ImageCollection}
  & Encapsula o estado da imagem (original e processada), parâmetros de
    ajuste (brilho, contraste, rotação, inversão) e a coleção de imagens
    abertas simultaneamente. \\
\addlinespace
View \\
\texttt{MainWindow}, \texttt{ImageViewWidget}
  & Gerencia a interface gráfica Qt, renderiza a imagem com suporte a
    zoom, pan e anotações, e exibe a lista lateral de imagens abertas. \\
\addlinespace
Controller \\
\texttt{ImageProcessor}
  & Contém algoritmos puros de processamento de imagem, completamente
    isolados da camada de interface. Facilita testes unitários
    independentes. \\
\bottomrule
\end{tabularx}
\end{table}

\subsection{Sistema de Build --- CMake}

Optou-se por \textbf{CMake} em detrimento do \texttt{qmake} por três razões
objetivas:

\begin{enumerate}
  \item \textbf{Suporte a IDEs modernas:} integração nativa com VSCode
        (\textit{CMake Tools}), CLion e Qt Creator 6+.
  \item \textbf{Separação de targets:} o projeto principal e os testes
        unitários são targets independentes, compilados e executados
        separadamente via \texttt{ctest}.
  \item \textbf{Adoção da indústria:} CMake é o padrão de facto em projetos
        C++ corporativos modernos.
\end{enumerate}

\subsection{Linguagem --- C++17}

O padrão C++17 foi escolhido por oferecer ferramentas que tornam o código
mais seguro e expressivo sem custo de runtime:

\begin{itemize}
  \item \texttt{std::optional<QImage>} para representar estado ``sem imagem
        carregada'' de forma type-safe, eliminando sentinelas frágeis como
        \texttt{QImage()} vazio ou flags booleanas separadas.
  \item \texttt{std::clamp} para garantir limites de valores sem risco de
        overflow silencioso.
  \item \texttt{std::array<uint8\_t, 256>} alocado em stack para a LUT,
        evitando alocação dinâmica em hot paths.
  \item \texttt{if constexpr}, structured bindings e deduções de tipo
        (\texttt{auto}) onde aplicável.
\end{itemize}

\subsection{Padronização e Nomenclatura}

Embora o presente relatório e a interface final do usuário utilizem o idioma
português para facilitar a comunicação, optou-se pela \textbf{nomenclatura em
inglês} para todos os elementos técnicos e de desenvolvimento do projeto,
incluindo:

\begin{itemize}
  \item Nomes de arquivos e diretórios da estrutura de código;
  \item Nomes de classes, funções, variáveis e membros;
  \item Mensagens de \textit{commit} no histórico do Git.
\end{itemize}

Esta decisão alinha o projeto às melhores práticas da indústria global de
\textit{software}, garantindo consistência com as bibliotecas padrão (STL), o
\textit{framework} Qt e facilitando a futura manutenção por equipes
multiculturais ou o uso de ferramentas internacionais de análise de código.

% ========================================================================
\section{Formatos de Imagem Suportados}

O enunciado solicitava suporte a pelo menos um dos formatos BMP, PNG ou JPEG
(\textit{bônus} para mais de um). O projeto suporta \textbf{os três formatos},
sem dependência de bibliotecas externas, pois o Qt 6 os gerencia nativamente
via \texttt{QImage::load()}.

\begin{table}[h!]
\centering
\caption{Formatos de imagem suportados}
\vspace{4pt}
\begin{tabular}{lll}
\toprule
Formato & Extensões & Característica relevante \\
\midrule
BMP  & \texttt{.bmp}              & Sem compressão, leitura trivial \\
PNG  & \texttt{.png}              & Canal alfa; convertido para RGB32 \\
JPEG & \texttt{.jpg}, \texttt{.jpeg} & Compressão com perda; mais comum em PACS \\
\bottomrule
\end{tabular}
\end{table}

Todas as imagens são convertidas imediatamente para o formato interno
\texttt{QImage::Format\_RGB32} no momento do carregamento. Isso garante que os
algoritmos de processamento operem sempre sobre um layout de memória uniforme
(4 bytes por pixel: R, G, B + padding), independente do formato original.

% ========================================================================
\section{Técnicas Aplicadas}

\subsection{Lookup Table (LUT) para Brilho e Contraste}

A técnica de \textbf{Lookup Table} é o método profissional para aplicar
transformações tonais em imagens. A ideia central é:

\begin{infobox}[title={Princípio da LUT}]
  Em vez de calcular a fórmula de brilho/contraste para cada um dos
  potencialmente milhões de pixels, calcula-se apenas \textbf{256 vezes}
  (uma por possível valor de canal 0--255) e armazena-se o resultado em uma
  tabela. Cada pixel é então transformado por uma simples consulta de índice:
  $O(1)$ por pixel.
\end{infobox}

A fórmula aplicada combina contraste e brilho em uma única passagem:

\begin{equation}
  \text{LUT}[i] = \text{clamp}\!\left(
    f_c \cdot (i - 128) + 128 + B,\; 0,\; 255
  \right)
\end{equation}

onde $f_c = \frac{C + 100}{100}$ é o fator de contraste (intervalo $[0, 2]$)
e $B \in [-100, 100]$ é o deslocamento de brilho. O \texttt{std::clamp} garante
que o resultado nunca ultrapasse os limites do tipo \texttt{uint8\_t}, evitando
overflow silencioso.

O acesso à memória da imagem é feito diretamente via \texttt{scanLine(y)},
evitando a sobrecarga de chamadas ao método \texttt{pixel(x, y)}:

\begin{lstlisting}[caption={Aplicação da LUT com acesso direto à memória}]
for (int y = 0; y < h; ++y) {
    QRgb* line = reinterpret_cast<QRgb*>(m_processed.scanLine(y));
    for (int x = 0; x < w; ++x) {
        const QRgb px = line[x];
        line[x] = qRgb(lut[qRed(px)],
                       lut[qGreen(px)],
                       lut[qBlue(px)]);
    }
}
\end{lstlisting}

\subsection{Renderização de Alta Precisão (Pixel-Perfect)}

Para apoiar o diagnóstico médico detalhado, o widget de visualização implementa uma lógica de renderização adaptativa baseada no nível de ampliação:

\begin{itemize}
  \item \textbf{Zoom < 10$\times$:} utiliza interpolação bilinear (\texttt{Qt::SmoothTransformation}) para suavizar transições e evitar aliasing.
  \item \textbf{Zoom > 10$\times$:} a suavização é desativada (\texttt{nearest-neighbor}), permitindo que o médico visualize os valores reais dos pixels sem interpolação artificial que poderia mascarar artefatos.
  \item \textbf{Zoom > 15$\times$:} ativa-se uma \textbf{grade de contorno de pixels}. Utilizando uma \texttt{QPen} cosmética (largura fixa de 1px independente da escala), o sistema desenha bordas translúcidas ao redor de cada pixel, facilitando a contagem e análise de estruturas em nível microscópico.
\end{itemize}

\subsection{Integração de Controles e MenuUI}

A interface foi reformulada para oferecer flexibilidade total ao usuário, inspirada em estações de trabalho radiológicas (\textit{Workstations}) profissionais:

\begin{itemize}
    \item \textbf{Agrupamento Funcional:} Os controles na barra lateral são organizados em \texttt{QGroupBox} (Ajustes Tonais, Zoom, Rotação).
    \item \textbf{Visibilidade Dinâmica:} O menu ``Ferramentas'' permite que o usuário ligue/desligue a exibição de cada grupo individualmente, limpando a interface de elementos desnecessários para o fluxo atual.
    \item \textbf{Persistência de Visão:} Adição do menu ``Exibir'', que permite reabrir a barra lateral de ajustes caso ela seja fechada acidentalmente, garantindo que o usuário nunca perca o acesso aos controles.
\end{itemize}

\subsection{Ferramenta Caliper (Régua Médica)}

A régua de medição (\textit{caliper}) é a funcionalidade de maior complexidade
geométrica do projeto. Ela exige o mapeamento bidirecional entre dois espaços
de coordenadas:

\begin{itemize}
  \item \textbf{Coordenadas do widget:} onde o usuário clica com o mouse.
  \item \textbf{Coordenadas lógicas da imagem:} posição real no pixel original,
        independente de zoom, pan ou rotação aplicados.
\end{itemize}

A conversão inversa requer desfazer as transformações acumuladas pelo
\texttt{QPainter} (translação de pan, escala de zoom, rotação), o que foi
resolvido através da \textbf{inversa da matriz de transformação} composta:

\begin{equation}
  \mathbf{p}_{\text{imagem}} = \mathbf{M}^{-1} \cdot \mathbf{p}_{\text{widget}}
\end{equation}

onde $\mathbf{M} = T_{\text{pan}} \cdot S_{\text{zoom}} \cdot R_{\theta}$.

A medição é exibida em \textbf{milímetros}, utilizando uma calibração simulada
de $0{,}1\,\text{mm/pixel}$ --- valor representativo de modalidades de imagem
médica como radiografia digital. Em um sistema PACS real, esse valor seria
lido diretamente dos metadados DICOM (\textit{Pixel Spacing}).

\subsection{Inversão de Polaridade (Negativo)}

O filtro de inversão ($\text{valor}' = 255 - \text{valor}$) simula a
visualização de filmes radiográficos tradicionais, onde ossos aparecem escuros
sobre fundo claro. É implementado via \texttt{QImage::invertPixels()} do Qt,
operação otimizada internamente pela biblioteca.

% ========================================================================
\section{Funcionalidades Implementadas}

\subsection{Requisitos Obrigatórios}

\begin{itemize}
  \item Abertura de imagens via \texttt{QFileDialog} com filtro por extensão.
  \item Suporte a BMP, PNG e JPEG.
  \item Exibição centralizada da imagem com fundo preto.
  \item Controle de brilho via \texttt{QSlider} (intervalo $[-100, 100]$).
  \item Controle de contraste via \texttt{QSlider} (intervalo $[-100, 100]$).
  \item Sliders desabilitados enquanto nenhuma imagem estiver aberta.
\end{itemize}

\subsection{Funcionalidades Bônus}

\begin{itemize}
  \item \textbf{Bônus 4c --- Drag do mouse:} arrastar horizontalmente altera
        o brilho; verticalmente, o contraste. Os sliders são atualizados
        em tempo real.
  \item \textbf{Bônus 5 --- Múltiplas imagens:} lista lateral com miniaturas;
        seleção troca a imagem exibida sem fechar as demais.
  \item \textbf{Bônus 6a --- Zoom:} scroll do mouse com interpolação bilinear e centralizando imagem com pan.
  \item \textbf{Bônus 6b --- Pan:} arrasto com botão direito; imagem retorna ao
        centro com duplo clique.
  \item \textbf{Bônus 6c --- Rotação:} controle de rotação em 90° e ângulo
        livre por slider.
  \item \textbf{Bônus 7 --- Surpresa:} Ferramenta Caliper (régua calibrada em
        mm) e inversão de polaridade radiográfica.
\end{itemize}

% ========================================================================
\section{Testes Unitários}

Os testes foram implementados com o framework \textbf{Qt Test (QTest)},
incluído nativamente no Qt 6, dispensando dependências externas. Cada caso de
teste é um método de slot privado da classe de teste:

\begin{table}[h!]
\centering
\caption{Casos de teste implementados}
\vspace{4pt}
\begin{tabularx}{\linewidth}{lX}
\toprule
Caso de teste & O que verifica \\
\midrule
\texttt{testNeutralAdjustment}
  & Brilho $= 0$, contraste $= 0$ retorna imagem bitwise-idêntica à original. \\
\texttt{testBrightnessClamp}
  & Brilho $= +100$ em pixel $(250, 250, 250)$ resulta em $(255, 255, 255)$;
    sem overflow. \\
\texttt{testNegativeBrightness}
  & Brilho $= -100$ em pixel $(10, 10, 10)$ resulta em $(0, 0, 0)$;
    sem underflow. \\
\texttt{testContrastMidpoint}
  & Pixel exato em $128$ permanece $128$ para qualquer fator de contraste. \\
\texttt{testLoadInvalidFile}
  & \texttt{loadFromFile()} retorna \texttt{false} e
    \texttt{isLoaded()} permanece \texttt{false} com caminho inválido. \\
\texttt{testLoadValidFile}
  & \texttt{loadFromFile()} retorna \texttt{true} e imagem tem dimensões
    corretas com arquivo PNG válido. \\
\bottomrule
\end{tabularx}
\end{table}

Os testes são executados via \texttt{ctest} a partir do diretório de build:

\begin{lstlisting}[language=bash, caption={Execução dos testes}]
cd build && ctest --output-on-failure
\end{lstlisting}

% ========================================================================
\section{Dificuldades Encontradas}

\subsection{LUT com Canal Separado e Overflow}

A principal dificuldade técnica foi garantir que a LUT tratasse cada canal RGB
de forma independente, sem que a transformação de contraste --- centralizada
em 128 e de natureza não-linear --- interagisse incorretamente com o
deslocamento de brilho. O risco concreto era um \textbf{overflow silencioso}
ao fazer \texttt{static\_cast<uint8\_t>} em valores fora do intervalo
$[0, 255]$, que em C++ é comportamento indefinido para inteiros com sinal. A
solução foi aplicar \texttt{std::clamp} sobre o tipo \texttt{double} antes da
conversão, garantindo sempre um valor seguro.

\subsection{Gestão de Coordenadas Multi-Transformação}

Converter a posição do clique do mouse (espaço do widget) para a coordenada
real na imagem original, após zoom, pan e rotação acumulados, exigiu a
\textbf{inversão da matriz de transformação composta} do \texttt{QPainter}. O
Qt disponibiliza \texttt{QTransform::inverted()}, mas a construção manual da
composição correta de transformações (ordem importa: rotação antes de escala
antes de translação) demandou atenção especial para a ferramenta caliper
funcionar com precisão.

\subsection{Interpolação Adaptativa no Zoom}

Inicialmente, a imagem era sempre escalada com \texttt{SmoothTransformation},
o que borrava os pixels ao ampliar. A solução foi alternar para
\texttt{FastTransformation} (nearest-neighbor) quando o zoom ultrapassa $1\times$,
e adicionar a grade de pixels a partir de $8\times$, replicando o comportamento
de editores profissionais como o GIMP.

% ========================================================================
\section{Possíveis Melhorias Futuras}

\begin{enumerate}
  \item \textbf{Processamento em GPU via OpenGL.} Migrar os filtros de
        brilho/contraste para \textit{shaders} GLSL utilizando
        \texttt{QOpenGLWidget}, eliminando a latência CPU para imagens de
        alta resolução (ex.: radiografias $4096 \times 4096$).

  \item \textbf{Suporte nativo a DICOM via DCMTK.} Integrar a biblioteca
        \href{https://dicom.offis.de/dcmtk}{DCMTK} para leitura de metadados
        reais de calibração (\textit{Pixel Spacing}), eliminando a calibração
        simulada atual e tornando as medições do caliper clinicamente válidas.

  \item \textbf{Janelamento (\textit{Windowing}).} Implementar controles de
        \textit{Window Level / Window Width} para imagens de 12-bit ou 16-bit,
        padrão em Tomografia Computadorizada e Ressonância Magnética.

  \item \textbf{Processamento assíncrono.} Mover a aplicação da LUT para uma
        thread separada via \texttt{QtConcurrent::run}, mantendo a interface
        responsiva em imagens grandes.

  \item \textbf{Filtros convolucionais.} Adicionar filtros como
        \textit{sharpening}, suavização gaussiana e detecção de bordas
        (Sobel/Canny) via \texttt{std::array} de kernels aplicados com
        \texttt{QImage}.
\end{enumerate}

% ========================================================================
\section{Avaliação do Resultado}

O projeto entregou todos os requisitos obrigatórios e todos os bônus
especificados no enunciado, acrescido da funcionalidade extra de caliper
médico. A arquitetura MVC adotada manteve o código testável e extensível: o
\texttt{ImageProcessor} não tem dependência alguma de Qt Widgets, podendo ser
testado em isolamento completo.

A decisão de implementar funcionalidades orientadas ao domínio médico ---
calibração em mm, inversão de polaridade, interpolação adaptativa para
diagnóstico --- foi deliberada: demonstra compreensão do contexto de negócio
da Pixeon (PACS, imagens médicas, DICOM) e da responsabilidade que esse
software carrega para profissionais de saúde.

Do ponto de vista pessoal, a maior evolução foi a gestão de coordenadas
multi-transformação para o caliper, que exigiu revisitar álgebra linear
aplicada --- um domínio fora da zona de conforto inicial, mas essencial para
ferramentas de medição precisas em ambiente clínico.

\begin{infobox}[title={Consideração Final}]
  Em um contexto de imagens médicas como DICOM, o pipeline de processamento
  implementado --- LUT de 256 entradas, acesso direto a \texttt{scanLine},
  interpolação adaptativa --- pode ser expandido para suportar imagens de
  16 bits e metadados DICOM reais, garantindo fidelidade diagnóstica sem
  comprometer a performance.
\end{infobox}

% --- Fim do documento ---
\end{document}